\documentclass[12pt]{article}

\usepackage[a4paper,margin=1in]{geometry}
\usepackage{amsmath, amssymb, amsthm}
\usepackage{enumitem}
\usepackage{titlesec}
\usepackage{hyperref}

\titleformat{\section}{\large\bfseries}{\thesection}{1em}{}
\titleformat{\subsection}{\normalsize\bfseries}{\thesubsection}{1em}{}

\title{\textbf{CSE400 -- Fundamentals of Probability in Computing}}
\author{Lecture 4: Joint Probability and Conditional Probability}
\date{Instructor: Dhaval Patel, PhD \\ January 15, 2026}

\begin{document}
\maketitle
\hrule
\vspace{1em}

\section{Course Context and Motivation}

\subsection{Why Learn CSE400?}

Daily life conversations involve uncertainty.

\subsection{Engineering Applications}

\begin{itemize}
\item Speech Recognition
\item Radar Systems
\item Communication Networks
\end{itemize}

\subsubsection*{Speech Recognition System}
Vocabulary templates such as \textbf{Hello, Yes, No, Bye} are matched against speech signals from different speakers.

\subsubsection*{Radar System}
Radar systems detect objects by transmitting signals and analyzing returned reflections.

\subsubsection*{Communication Network}
Information travels from a \textbf{source} to a \textbf{destination} through interconnected devices.

\section{Lecture Outline}

\subsection*{Introduction to Probability Theory}
\begin{itemize}
\item Experiments, Sample Space, and Events
\item Axioms of Probability
\item Corollaries and Propositions from Probability Axioms
\item Assigning Probability: Classical and Relative Frequency Approaches
\end{itemize}

\subsection*{Joint Probability}
\begin{itemize}
\item Motivation, notation, and concepts
\item Example: Card deck
\item Example: Costume party
\end{itemize}

\subsection*{Conditional Probability}
\begin{itemize}
\item Motivation, notation, and concepts
\item Cards without replacement
\item Game of poker
\item The missing key
\end{itemize}

\textit{(Slides provided cover the Introduction to Probability Theory section.)}

\section{Introduction to Probability Theory}

\subsection{Experiments, Outcomes, and Events}

\subsubsection*{Experiment ($E$)}
A procedure performed that produces a result.

\textbf{Example:} Tossing a coin five times ($E_5$).

\subsubsection*{Outcome ($\xi$)}
A possible result of an experiment.

\textbf{Example:}
\[
\xi_1 = HHTHT
\]

\subsubsection*{Event}
A set of outcomes of an experiment.

\textbf{Example:} Event $C$ with experiment $E_5$:
\[
C = \{\text{all outcomes consisting of an even number of heads}\}
\]

\subsection{Sample Space}

\subsubsection*{Sample Space ($S$)}
The set of all possible distinct outcomes of an experiment.

\subsubsection*{Properties}

\begin{itemize}
\item \textbf{Mutually Exclusive:} heads or tails, but not both.
\item \textbf{Collectively Exhaustive:} no outcome other than heads or tails.
\end{itemize}

\subsubsection*{Additional Notes}

\begin{itemize}
\item $S$ is the universal set of outcomes.
\item Sample space may be:
\begin{itemize}
\item Discrete
\item Countably infinite
\item Continuous
\end{itemize}
\end{itemize}

\subsubsection*{Examples of Sample Spaces}

\begin{enumerate}
\item Flipping a fair coin once
\item Rolling a cubical die
\item Rolling two dice
\item Flip a coin until a tails occurs
\item Random number generator on interval $[0,1]$
\end{enumerate}

\section{Probability and Axioms}

\subsection{Probability}

Probability measures the likelihood of events, or a function assigning a numerical value representing likelihood.

\subsection{Axioms of Probability}

These statements are self-evident and require no proof.

\subsubsection*{Axiom 1 (Bounds)}

\[
0 \le Pr(A) \le 1
\]

\subsubsection*{Axiom 2 (Certainty)}

If $S$ is the sample space:
\[
Pr(S) = 1
\]

\subsubsection*{Axiom 3 (Additivity)}

If $A \cap B = \emptyset$:
\[
Pr(A \cup B) = Pr(A) + Pr(B)
\]

For infinitely many mutually exclusive events $A_i$:

\[
Pr\left(\bigcup_{i=1}^{\infty} A_i\right)
= \sum_{i=1}^{\infty} Pr(A_i)
\]

\section{Corollary from Probability Axioms}

Define $A_i$ as the null event for all values of $i>n$.

\subsection*{Corollary 2.1}

For $M$ finite mutually exclusive sets $A_i$:

\[
Pr\left(\bigcup_{i=1}^{M} A_i\right)
= \sum_{i=1}^{M} Pr(A_i)
\]

\textbf{Notes}
\begin{itemize}
\item Axiom 3 is equivalent to Corollary 2.1 for finite sample spaces.
\item Axiom 3 is needed when the sample space is infinite.
\end{itemize}

\section{Propositions from Probability Axioms}

\subsection*{Proposition 2.1 (Complement Rule)}

\[
Pr(A^c) = 1 - Pr(A)
\]

\subsection*{Proposition 2.2 (Monotonicity)}

If $A \subset B$, then:
\[
Pr(A) \le Pr(B)
\]

\section*{Key Points for Exam Preparation}

\begin{itemize}
\item Experiment $\rightarrow$ procedure producing results.
\item Outcome $\rightarrow$ possible result.
\item Event $\rightarrow$ set of outcomes.
\item Sample space $\rightarrow$ all possible outcomes.
\item Outcomes must be mutually exclusive and collectively exhaustive.
\item Probability follows three axioms.
\item Additivity applies to mutually exclusive events.
\item Complement and subset properties follow from axioms.
\item Finite additivity is given by Corollary 2.1.
\end{itemize}

\vfill
\hrule
\centerline{\textit{End of Lecture 4 Scribe}}

\end{document}
